\documentclass[12pt]{article}
\usepackage[utf8]{inputenc}
\usepackage[spanish]{babel}
\usepackage{amsmath}
\usepackage{amsfonts}
\usepackage{amssymb}
\usepackage{geometry}
\geometry{letterpaper, margin=2.5cm}

\begin{document}

\begin{center}
    \Large \textbf{Evaluación de Examen 1: Unidad I} \\
    \large Física para Ciencias de la Salud (005-1713) \\
    \vspace{0.5cm}
    \textbf{Estudiante:} Franyerika Rodriguez \quad \textbf{C.I.:} 33.278.284
    \vspace{0.3cm} \\
    \large \textbf{Calificación Final: 8/10 puntos}
\end{center}

\vspace{0.5cm}

\section*{Análisis de Resultados}

\subsection*{1. Ángulo plano (Conversión a radianes)}
\textbf{Procedimiento y resultado:} La estudiante demuestra gran seguridad al realizar la regla de tres mediante dos equivalencias distintas ($90^\circ$ y $180^\circ$). En ambos casos el desarrollo aritmético es impecable, llegando al resultado correcto de $0,41\pi$ rad (aproximación de $\frac{5\pi}{12}$). \\
\textbf{Veredicto:} \textbf{Correcto} (2/2 pts).

\subsection*{2. Sistemas de coordenadas (Longitud del hueso)}
\textbf{Procedimiento y resultado:} Excelente resolución geométrica. Dibuja el gráfico, identifica con precisión los catetos del triángulo rectángulo formado ($a=6$, $b=8$) y aplica el Teorema de Pitágoras ($\sqrt{36+64}=\sqrt{100}$) para obtener $L = 10$ cm. \\
\textbf{Veredicto:} \textbf{Correcto} (2/2 pts).

\subsection*{3. Vectores (Fuerza resultante)}
\textbf{Procedimiento y resultado:} Plantea correctamente la suma de las componentes $\hat{i}$ y $\hat{j}$. Ejecuta la resta de $45 - 15 = 30$ y la suma de $25 + 35 = 60$ respetando los signos, para obtener el vector resultante exacto de $30\hat{i} + 60\hat{j}$ N. \\
\textbf{Veredicto:} \textbf{Correcto} (2/2 pts).

\subsection*{4. Potencias de diez (Notación científica y prefijo)}
\textbf{Procedimiento y resultado:} Logra expresar el número en notación científica ($1,25 \times 10^{-5}$) y lo ajusta a $10^{-6}$ con la intención de colocar el prefijo. No obstante, utiliza el símbolo "M" (que corresponde a Mega, $10^6$) en lugar del símbolo correcto "$\mu$" (micro, $10^{-6}$), escribiendo $12,5$ Mm en vez de $12,5\,\mu$m. \\
\textbf{Veredicto:} \textbf{Parcialmente correcto} (Desarrollo numérico correcto, error en la simbología del S.I.) (1/2 pts).

\subsection*{5. Sistemas de unidades (Conversión de densidad)}
\textbf{Procedimiento y resultado:} Decide resolver el problema haciendo conversiones separadas por reglas de tres. La conversión de la masa es correcta ($1,03 \times 10^{-3}$ kg). Sin embargo, comete un error algebraico al despejar la regla de tres del volumen, multiplicando por el factor $10^6$ en lugar de dividir, obteniendo un volumen erróneo que afecta toda la operación final ($1,03 \times 10^{-9} \text{ kg/m}^3$). \\
\textbf{Veredicto:} \textbf{Parcialmente correcto} (Conversión de masa correcta, error en el despeje del volumen) (1/2 pts).

\vspace{0.5cm}
\section*{Comentarios Generales y Sugerencias}
La estudiante evidencia un pensamiento muy lógico, con un gran nivel de detalle y validación en sus respuestas (como se observa en la Pregunta 1). 
\begin{itemize}
    \item Se recomienda repasar los símbolos de los prefijos del Sistema Internacional (diferenciar $\mu$ de M).
    \item Para conversiones de unidades derivadas (como la densidad), suele ser menos propenso a errores utilizar el método de "factores de conversión en cadena" en lugar de hacer reglas de tres separadas.
\end{itemize}

\end{document}